\chapter{The Anatomy of an Agent}

An agentic AI system combines a foundation model with the infrastructure required to
make it useful: tools, memory, security. This book is organised around a decomposition of Agentic AI into four layers: the intelligence layer, the action layer, the memory and
knowledge layer, and the security layer. These layers are conceptually distinct but
continuously interdependent in production: the model reasons over context assembled by
the memory layer, selects actions exposed by the action layer, and operates within
boundaries enforced by the security layer. Evaluation and observability sit across all
four, measuring whether the full system behaves correctly at each step and over time.

\section{The Four Layers of an Agent}

The first layer is the \emph{intelligence layer}. In this book, the intelligence layer
encompasses two tightly coupled components: the foundation model, which provides the
core reasoning and generation capability, and the cognitive architecture, which governs
how the model pursues a goal across multiple steps, deciding when to act, when to
retrieve, when to seek clarification, and when to stop. The subsequent chapter explains
how these models work, beginning with the original Transformer architecture and
progressing toward the design choices that characterise modern frontier models.

The second layer is the \emph{action layer}. This layer enables the agent to produce
effects beyond generating text. It encompasses function calling, APIs, tools, MCP
servers and agent-as-tool patterns. Some tools are
deterministic, such as a calculator, while others are
stochastic, such as a sub-agent that performs research, writes code, or produces a
judgment. Together they define the set of effects the agent can produce in the world
beyond generating text.

The third layer is the \emph{memory and knowledge layer}. This layer determines what
information is available to the model at the moment of decision. It includes short-term
conversation state, long-term user and organisational memory, retrieval-augmented
generation, knowledge bases, documents, vector stores and reusable
skills. In this book, skills are treated as a form of procedural knowledge: versioned,
retrievable packages of instructions or patterns that improve the reliability with which
the agent performs recurring tasks. The quality of this layer is as consequential as the
capability of the model that draws on it: an agent operating from incomplete, stale, or
poorly structured context will produce unreliable outputs regardless of the underlying
model's capacity.

The fourth layer is the \emph{security layer}. An agent with access to tools, persistent
memory, and the ability to take consequential actions requires security as a foundational
architectural concern. The security layer includes identity, authentication,
authorisation, permission boundaries, guardrails, human approval mechanisms, audit
logging, sandboxing, and protection against prompt injection and data exfiltration. The
scope and rigour of this layer should scale with the autonomy and consequence of the
agent's permitted actions.

\section{From Anatomy to Implementation}

The architecture described above provides the conceptual foundation for the rest of the
book. The chapters that follow develop each layer in depth: how foundation models are
trained and how they reason, how the action layer is designed and governed, how the
memory and knowledge layer is constructed and evaluated, and how the security layer is
enforced in production. The application section then demonstrates how these layers
compose into an end-to-end system: an agent built with LangGraph, connected to tools
and memory, evaluated with Langfuse, and iterated through a pipeline that closes the
loop between offline evaluation, production traces, human feedback, and future evaluation
datasets. The goal is to give the reader the conceptual and practical tools to build an
agent that can be tested, monitored, improved, and deployed through a governed production
pipeline in their own domain.
